\documentclass[10pt,compress,aspectratio=169]{beamer}
\usepackage[utf8]{inputenc}
\usepackage[OT1]{fontenc}
\usepackage{amsmath,amsthm,amssymb}
\usepackage{mathpartir}
\usepackage{stmaryrd}
\usepackage{booktabs}
\usepackage{minted}
\setminted{fontsize=\footnotesize, breaklines=true, autogobble=true}

% Load Theme
\usetheme[
	navigation=false,
	frametotal=false,
	maincolor=rptudunkelblau,
	oneColorBoxes,
	signetontitle=true,
	redhat=false
]{RPTU}

\title{Deriving Verified Abstract Machines}
\subtitle{From Big-Step Semantics in Agda}
\date[May 2026]{May 2026}
\author[M.\,Hamido]{Mahmoud M. Hamido}
\institute[]{Department of Computer Science\\RPTU Kaiserslautern-Landau}

% Convenience macro: highlight where we are in the equivalence chain
\newcommand{\chainhere}[1]{%
  \begin{center}\small
  $\text{Denotational}
  \;\leftrightarrow\;
  \text{Big-Step}
  \;\leftrightarrow\;
  \mathbf{#1}
  \;\leftrightarrow\;
  \text{Small-Step}$
  \end{center}}

\begin{document}
\maketitle

\begin{frame}{Outline}
	\tableofcontents
\end{frame}

%%%%%%%%%%%%%%%%%%%%%%%%%%%%%%%%%%%%%%%%%%%%%%%%%%%%%%%%%%%%%%%%%%%%%%%%%%%
\section{Motivation}
%%%%%%%%%%%%%%%%%%%%%%%%%%%%%%%%%%%%%%%%%%%%%%%%%%%%%%%%%%%%%%%%%%%%%%%%%%%

\begin{frame}[coloredbackground,stackedUbackground]{}
  \centering\vfill
  {\usebeamerfont{title}\usebeamercolor[fg]{title}Part I\par}
  \bigskip
  {\Large\bfseries Motivation}
  \vfill
\end{frame}

\begin{frame}{The Problem: Specs vs.\ Implementations}
  \begin{columns}[T]
    \column{0.52\textwidth}
      \begin{block}{The Expectation}
        A language implementation should faithfully realise its specification.
      \end{block}
      \bigskip
      In practice, specifications are \textbf{written documents}:
      \begin{itemize}
        \item C++ Standard
        \item Java Language Specification
        \item WebAssembly Specification
      \end{itemize}
      \bigskip
      Conformance is checked by testing against reference implementations.
    \column{0.44\textwidth}
      \begin{alertblock}{The Fundamental Flaw}
        Testing covers \textbf{finitely many} programs.\\[6pt]
        Specifications make claims about \textbf{all} programs.\\[6pt]
        No matter how thorough the tests, some combination of features may expose a discrepancy.
      \end{alertblock}
  \end{columns}
\end{frame}

\begin{frame}{Formal Verification \& This Thesis}
  \begin{block}{The Solution}
    \textbf{Prove by construction} that an implementation conforms to the specification, covering \emph{all} inputs simultaneously.
  \end{block}
  \bigskip
  \begin{columns}[T]
    \column{0.48\textwidth}
      \textbf{Tool:} Agda --- a dependently typed proof assistant\\[8pt]
      \textbf{Foundation:} Curry--Howard correspondence
      \begin{align*}
        \text{proposition} &\;\longleftrightarrow\; \text{type}\\
        \text{proof}       &\;\longleftrightarrow\; \text{value}
      \end{align*}
      The type checker \emph{is} the proof checker.
    \column{0.48\textwidth}
      \begin{block}{Central Thesis}
        An implementation can be \emph{derived directly} from its specification.
      \end{block}
      \bigskip
      \textbf{Three case studies:}
      \begin{enumerate}
        \item Arithmetic expressions
        \item Arithmetic with exceptions
        \item Simply Typed Lambda Calculus
      \end{enumerate}
  \end{columns}
\end{frame}

%%%%%%%%%%%%%%%%%%%%%%%%%%%%%%%%%%%%%%%%%%%%%%%%%%%%%%%%%%%%%%%%%%%%%%%%%%%
\section{Background}
%%%%%%%%%%%%%%%%%%%%%%%%%%%%%%%%%%%%%%%%%%%%%%%%%%%%%%%%%%%%%%%%%%%%%%%%%%%

\begin{frame}[coloredbackground,stackedUbackground]{}
  \centering\vfill
  {\usebeamerfont{title}\usebeamercolor[fg]{title}Part II\par}
  \bigskip
  {\Large\bfseries Background}
  \vfill
\end{frame}

\begin{frame}[fragile]{Agda: Dependent Types and Propositions as Types}
  Types may \textbf{depend on values} --- invariants are encoded in types.
  \bigskip
  \begin{columns}[T]
    \column{0.48\textwidth}
      \begin{minted}[bgcolor=gray!8]{agda}
data Vec (A : Set) : N -> Set where
  []  : Vec A zero
  _::_ : A -> Vec A n -> Vec A (suc n)
      \end{minted}
      {\footnotesize Length tracked at the type level.\\Ill-typed terms are not representable.}
    \column{0.48\textwidth}
      \begin{minted}[bgcolor=gray!8]{agda}
data _==_ {A : Set}(x : A) : A -> Set where
  refl : x == x

-- A well-typed program is a valid proof
+-comm : forall n m -> n + m == m + n
      \end{minted}
      {\footnotesize Agda checks termination to prevent circular proofs.}
  \end{columns}
\end{frame}

\begin{frame}[fragile]{Intrinsic Typing}
  Expressions are \textbf{indexed by their type}. Every constructible term is automatically well-typed.
  \bigskip
  \begin{columns}[T]
    \column{0.48\textwidth}
      \begin{minted}[bgcolor=gray!8]{agda}
data Type : Set where
  Nat  : Type
  Bool : Type

-- Value domain
|= Nat  = N
|= Bool = B
      \end{minted}
    \column{0.48\textwidth}
      \begin{minted}[bgcolor=gray!8]{agda}
-- Expressions indexed by Type
data |-_ : Type -> Set where
  `_            : |= T -> |- T
  _+_           : |- Nat -> |- Nat
                -> |- Nat
  if_then_else_ : |- Bool -> |- T
                -> |- T -> |- T
      \end{minted}
      {\footnotesize The condition \emph{must} be a Bool expression.\\Both branches \emph{must} share the same type.}
  \end{columns}
\end{frame}

\begin{frame}{The Four Semantics}
  \begin{center}
  \begin{tabular}{llll}
    \toprule
    \textbf{Style} & \textbf{Describes} & \textbf{Notation} & \textbf{Character} \\
    \midrule
    Big-step (natural)       & $e$ evaluates to $v$ directly    & $e \Downarrow v$              & Concise \\[4pt]
    Small-step (operational) & One-step reductions               & $e \longrightarrow e'$        & Verbose, fine-grained \\[4pt]
    Denotational             & Maps terms to a meaning domain   & $\llbracket e \rrbracket = v$ & Interpreter-like \\[4pt]
    Abstract machine         & Explicit states + transitions    & $s \Rightarrow s'$            & Operational \\
    \bottomrule
  \end{tabular}
  \end{center}
  \bigskip
  \begin{block}{Key Property}
    Under consistent semantics, they can be \textbf{related} --- and in some cases one can be \emph{derived} from another.
  \end{block}
\end{frame}

\begin{frame}{Proof Strategy: Chain of Equivalences}
  Rather than proving each pair of semantics equivalent in isolation:
  \[
    \text{Denotational}
    \;\longleftrightarrow\;
    \text{Big-Step}
    \;\longleftrightarrow\;
    \text{Machine}
    \;\longleftrightarrow\;
    \text{Small-Step}
  \]
  \begin{itemize}
    \item Each link is proved \textbf{independently}
    \item \textbf{Transitivity} gives equivalence between any two points
    \item An indirect equivalence can be composed from intermediate links
  \end{itemize}
  \bigskip
  \begin{block}{Analogy: Intermediate Representations in a Compiler}
    Rather than translating source directly to assembly in one step,
    incremental transformations via intermediate representations are used.
  \end{block}
  \bigskip
  {\small\textbf{Central focus:} the link $\text{Big-Step} \longleftrightarrow \text{Machine}$}
\end{frame}

%%%%%%%%%%%%%%%%%%%%%%%%%%%%%%%%%%%%%%%%%%%%%%%%%%%%%%%%%%%%%%%%%%%%%%%%%%%
\section{Case Study I: Arithmetic}
%%%%%%%%%%%%%%%%%%%%%%%%%%%%%%%%%%%%%%%%%%%%%%%%%%%%%%%%%%%%%%%%%%%%%%%%%%%

\begin{frame}[coloredbackground,stackedUbackground]{}
  \centering\vfill
  {\usebeamerfont{title}\usebeamercolor[fg]{title}Part III\par}
  \bigskip
  {\Large\bfseries Case Study I: Arithmetic Expressions}
  \vfill
\end{frame}

\begin{frame}[fragile]{Big-Step Semantics}
  \chainhere{Machine}
  \begin{columns}[T]
    \column{0.52\textwidth}
      \begin{minted}[bgcolor=gray!8]{agda}
data _|>_ : |- T -> |= T -> Set where
  `_      : (v : |= T) -> (` v) |> v

  _+_     : e1 |> n1 -> e2 |> n2
          -> (e1 + e2) |> (n1 + n2)

  if-then : e |> true  -> e1 |> v
          -> (if e then e1 else e2) |> v

  if-else : e |> false -> e2 |> v
          -> (if e then e1 else e2) |> v
      \end{minted}
    \column{0.44\textwidth}
      \begin{block}{Key Properties}
        \textbf{Deterministic:}\\[3pt]
        $e \Downarrow v_1 \;\land\; e \Downarrow v_2$\\[2pt]
        $\;\Rightarrow\; v_1 = v_2$\\[8pt]
        \textbf{Total:}\\[3pt]
        $\forall e.\; \exists v.\; e \Downarrow v$
      \end{block}
      \bigskip
      Both proved by structural induction on the derivation / expression.
  \end{columns}
\end{frame}

\begin{frame}{From Interpreter to Abstract Machine}
  Consider an interpreter as a recursive function traversing the expression tree.
  \begin{block}{Key Idea: Reify Suspended Computations as \emph{Frames}}
    For $e_1 \oplus e_2$, the interpreter must:
    \begin{enumerate}
      \item Evaluate $e_1$ --- the whole expression is \emph{suspended} while waiting\\[2pt]
            $\Rightarrow$ push frame $\;(\circ \oplus e_2)$
      \item Once $v_1$ is known, the frame becomes $(v_1 \oplus \circ)$, awaiting $e_2$
    \end{enumerate}
  \end{block}
  \bigskip
  \begin{columns}[T]
    \column{0.5\textwidth}
      The collection of outstanding frames forms a \textbf{stack}.\\[4pt]
      The machine halts when it reaches a return state with an \textbf{empty stack}.
    \column{0.46\textwidth}
      Two machine modes:
      \begin{description}
        \item[\textbf{Eval} $\uparrow e$] evaluating expression $e$
        \item[\textbf{Return} $\downarrow v$] returning value $v$
      \end{description}
  \end{columns}
\end{frame}

\begin{frame}[fragile]{Machine: Frames, States, Transitions}
  \begin{columns}[T]
    \column{0.52\textwidth}
      \begin{minted}[bgcolor=gray!8]{agda}
data Frame : Type -> Type -> Set where
  _+0_ : Hole -> |- Nat -> Frame Nat Nat
  _+1_ : |= Nat -> Hole -> Frame Nat Nat
  _else_ : |- T -> |- T -> Frame T Bool

data State (A : Type) : Set where
  eval   : Stack A T -> |- T -> State A
  return : Stack A T -> |= T -> State A
      \end{minted}
    \column{0.44\textwidth}
      Selected transitions:
      \begin{minted}[bgcolor=gray!8]{agda}
-- Decompose addition
step-+1 :
  eval stk (e1 + e2)
  ~> eval (stk :: (+0 e2)) e1

-- Left known; evaluate right
step-+2 :
  return (stk :: (+0 e2)) v1
  ~> eval (stk :: (v1 +1)) e2

-- Both known; sum and return
step-+3 :
  return (stk :: (v1 +1)) v2
  ~> return stk (v1 + v2)
      \end{minted}
  \end{columns}
\end{frame}

\begin{frame}[fragile]{Correctness and Completeness}
  \begin{block}{Completeness: big-step $\Rightarrow$ machine}
    \begin{minted}[bgcolor=gray!8]{agda}
complete : e |> v -> (eval stk e) ~>* (return stk v)
    \end{minted}
    \begin{itemize}
      \item Proved by \textbf{induction on the big-step derivation}
      \item Must generalise over an \emph{arbitrary} stack
    \end{itemize}
  \end{block}
  \medskip
  \begin{block}{Correctness: machine $\Rightarrow$ big-step (via confluence)}
    \begin{minted}[bgcolor=gray!8]{agda}
correct : (eval [] e) ~>* (return [] v) -> e |> v
    \end{minted}
    \begin{enumerate}
      \item \textbf{Totality} gives a canonical big-step derivation to some $v'$
      \item \textbf{Completeness} simulates it into a machine run ending at $\downarrow v'$
      \item \textbf{Confluence}: two runs from $e$ to values must agree $\Rightarrow v = v'$
    \end{enumerate}
  \end{block}
  \smallskip
  {\footnotesize This confluence argument is \textbf{reused in all three case studies}.}
\end{frame}

%%%%%%%%%%%%%%%%%%%%%%%%%%%%%%%%%%%%%%%%%%%%%%%%%%%%%%%%%%%%%%%%%%%%%%%%%%%
\section{Case Study II: Exceptions}
%%%%%%%%%%%%%%%%%%%%%%%%%%%%%%%%%%%%%%%%%%%%%%%%%%%%%%%%%%%%%%%%%%%%%%%%%%%

\begin{frame}[coloredbackground,stackedUbackground]{}
  \centering\vfill
  {\usebeamerfont{title}\usebeamercolor[fg]{title}Part IV\par}
  \bigskip
  {\Large\bfseries Case Study II: Arithmetic with Exceptions}
  \vfill
\end{frame}

\begin{frame}[fragile]{Extending the Language with Exceptions}
  \begin{columns}[T]
    \column{0.52\textwidth}
      \begin{minted}[bgcolor=gray!8]{agda}
data Exception : Set where
  err : Exception

-- New expression forms
throw      : Exception -> |- T
try_catch_ : |- T
           -> (Exception -> |- T)
           -> |- T
      \end{minted}
      {\footnotesize A single \texttt{err} constructor suffices.\\A tagged exception type allows handlers to re-raise on non-matching tags.}
    \column{0.44\textwidth}
      Terms now evaluate to a \textbf{result}:
      \begin{minted}[bgcolor=gray!8]{agda}
data Result (T : Type) : Set where
  return : |= T       -> Result T
  raise  : Exception -> Result T
      \end{minted}
      \bigskip
      Convenient aliases:
      \begin{itemize}
        \item $e \Downarrow v \;\triangleq\; e \Downarrow^r \operatorname{return}\,v$
        \item $e \Uparrow ex \;\triangleq\; e \Downarrow^r \operatorname{raise}\,ex$
      \end{itemize}
  \end{columns}
\end{frame}

\begin{frame}[fragile]{Big-Step Semantics with Exceptions}
  \chainhere{Machine}
  Uniform pattern: one rule for \emph{normal return}, one for \emph{exception propagation}.
  \medskip
  \begin{minted}[bgcolor=gray!8]{agda}
data _|>r_ : |- T -> Result T -> Set where
  _+_     : e1 |> n1 -> e2 |> n2 -> (e1 + e2) |> (n1 + n2)
  +-exn1  : e1 ^ ex               -> (e1 + e2) ^ ex
  +-exn2  : e1 |> n1 -> e2 ^ ex  -> (e1 + e2) ^ ex

  throw-rule : throw ex |>r raise ex

  try-return : e |> v               -> (try e catch h) |> v
  try-catch  : e ^ ex -> h ex |>r r -> (try e catch h) |>r r
  \end{minted}
  The handler $h$ receives the exception and may itself succeed or raise further.
\end{frame}

\begin{frame}[fragile]{The Machine: Three Modes and Unwind}
  \begin{columns}[T]
    \column{0.52\textwidth}
      A new \textbf{unwind mode} carries an exception through the stack:
      \begin{minted}[bgcolor=gray!8]{agda}
data State (A : Type) : Set where
  eval   : Stack A T -> |- T      -> State A
  return : Stack A T -> |= T      -> State A
  unwind : Stack A T -> Exception -> State A
      \end{minted}
      New frame and key transitions:
      \begin{minted}[bgcolor=gray!8]{agda}
catch : (Exception -> |- T) -> Frame T T

step-throw :
  eval stk (throw ex) ~> unwind stk ex

step-catch-exn :
  unwind (stk :: catch h) ex ~> eval stk (h ex)
      \end{minted}
    \column{0.44\textwidth}
      \begin{block}{Per-frame drop rules}
        Each frame needs an \emph{explicit} drop rule when unwinding:
      \end{block}
      \begin{minted}[bgcolor=gray!8]{agda}
step-drop-else :
  unwind (stk :: (e1 else e2)) ex
  ~> unwind stk ex

step-drop-+0 :
  unwind (stk :: (+0 e2)) ex
  ~> unwind stk ex
      \end{minted}
      {\footnotesize A single generic drop rule makes the completeness proof impossible: without inspecting the frame we cannot determine which big-step constructor applies.}
  \end{columns}
\end{frame}

\begin{frame}[fragile]{Mutual Recursion in Completeness}
  \begin{minted}[bgcolor=gray!8]{agda}
complete     : e |> v  -> (eval stk e) ~>* (return stk v)
complete-exn : e ^ ex  -> (eval stk e) ~>* (unwind stk ex)
  \end{minted}
  \bigskip
  \begin{block}{Why mutual recursion?}
    \begin{itemize}
      \item The \texttt{try-catch} case of \texttt{complete} needs \texttt{complete-exn} to handle the body (which raised)
      \item \texttt{complete-exn} calls \texttt{complete} to handle the handler (which produces a value)
    \end{itemize}
  \end{block}
  \bigskip
  \textbf{Correctness} extended to two terminal states.\\[4pt]
  If the machine reaches the \emph{wrong} terminal state ($\downarrow$ vs.\ unwind), an absurd equality between distinct constructors is discharged by \texttt{ex-falso}.
\end{frame}

%%%%%%%%%%%%%%%%%%%%%%%%%%%%%%%%%%%%%%%%%%%%%%%%%%%%%%%%%%%%%%%%%%%%%%%%%%%
\section{Case Study III: STLC}
%%%%%%%%%%%%%%%%%%%%%%%%%%%%%%%%%%%%%%%%%%%%%%%%%%%%%%%%%%%%%%%%%%%%%%%%%%%

\begin{frame}[coloredbackground,stackedUbackground]{}
  \centering\vfill
  {\usebeamerfont{title}\usebeamercolor[fg]{title}Part V\par}
  \bigskip
  {\Large\bfseries Case Study III: Simply Typed Lambda Calculus}
  \vfill
\end{frame}

\begin{frame}[fragile]{Syntax and De Bruijn Indices}
  \begin{columns}[T]
    \column{0.52\textwidth}
      \begin{minted}[bgcolor=gray!8]{agda}
data _|-_ (G : Context) : Type -> Set where
  Nat  : N -> G |- Nat
  Bool : B -> G |- Bool
  _+_  : G |- Nat -> G |- Nat -> G |- Nat
  Var  : T in G -> G |- T
  lam  : (G , T1) |- T2 -> G |- (T1 => T2)
  _._  : G |- (T1 => T2) -> G |- T1 -> G |- T2
      \end{minted}
    \column{0.44\textwidth}
      \begin{block}{De Bruijn indices}
        Variables as \emph{membership proofs} in the typing context:
      \end{block}
      \begin{minted}[bgcolor=gray!8]{agda}
data _in_ : Type -> Context -> Set where
  Z : T in (G , T)
  S : T in G -> T in (G , U)
      \end{minted}
      \medskip
      \begin{description}
        \item[Identity $\lambda x.x$] $\lambda\, Z$
        \item[Constant $\lambda x.\lambda y.x$] $\lambda\,\lambda\, (S\,Z)$
      \end{description}
      {\footnotesize Scope safety is guaranteed by the type.}
  \end{columns}
\end{frame}

\begin{frame}[fragile]{Closures and Big-Step Semantics}
  \chainhere{Machine}
  $\lambda$-abstractions produce \textbf{closures}: a saved environment paired with the function body.
  \medskip
  \begin{columns}[T]
    \column{0.52\textwidth}
      \begin{minted}[bgcolor=gray!8]{agda}
record Closure (A B : Type) : Set where
  inductive
  constructor <_,_>
  field
    {Gc} : Context
    dc   : Environment Gc
    ec   : (Gc , A) |- B

|= (A => B) = Closure A B  -- mutually recursive
      \end{minted}
    \column{0.44\textwidth}
      \begin{minted}[bgcolor=gray!8]{agda}
data _|-_|>_ :
  Env G -> G |- T -> |= T -> Set where

  VAR : d |- Var x |> lookup d x

  ABS : d |- lam e |> < d , e >

  app : d |- e0 |> < dc , ec >
      -> d |- e1 |> v1
      -> (dc , v1) |- ec |> v
      -> d |- e0 . e1 |> v
      \end{minted}
  \end{columns}
  {\footnotesize Totality uses \texttt{\{-\# TERMINATING \#-\}}: the recursive call on the closure body is not structurally smaller.}
\end{frame}

\begin{frame}[fragile]{Denotational Semantics: Two Approaches}
  \begin{columns}[T]
    \column{0.48\textwidth}
      \begin{block}{Direct Semantics}
        $\llbracket A \Rightarrow B \rrbracket \;=\; \llbracket A \rrbracket \to \llbracket B \rrbracket$
      \end{block}
      \begin{minted}[bgcolor=gray!8]{agda}
interpret d (lam e) =
  \z -> interpret (d , z) e

interpret d (e1 . e2) =
  (interpret d e1)(interpret d e2)
      \end{minted}
      {\footnotesize Function types are Agda functions.}
    \column{0.48\textwidth}
      \begin{block}{Closure Semantics}
        $\llbracket A \Rightarrow B \rrbracket \;=\; \mathsf{Closure}\,A\,B$
      \end{block}
      \begin{minted}[bgcolor=gray!8]{agda}
evaluate d (lam e)    = < d , e >
evaluate d (e1 . e2)  =
  apply (evaluate d e1)(evaluate d e2)

apply < dc , ec > v = evaluate (dc , v) ec
      \end{minted}
      {\footnotesize Matches the big-step value domain exactly.}
  \end{columns}
  \bigskip
  The two are related by \textbf{logical relations}: $v \sim_T u$ for each type $T$.\\[3pt]
  For function types: applying related functions to related arguments produces related results.
\end{frame}

\begin{frame}[fragile]{STLC Abstract Machine: Environment in State}
  The environment is \textbf{carried in the machine state}: $\delta \vdash \mathtt{stk} \uparrow e$.
  \medskip
  \begin{columns}[T]
    \column{0.52\textwidth}
      Three frames for each phase of function application:
      \begin{minted}[bgcolor=gray!8]{agda}
data Frame : Type -> Type -> Set where
  app1 : Hole -> G |- A   -> Frame B (A => B)
  app2 : |= (A => B) -> Hole -> Frame B A
  app3 : Env G -> |= (A => B) -> |= A
       -> Hole -> Frame B B
      \end{minted}
    \column{0.44\textwidth}
      \begin{block}{Three phases of $e_1 \cdot e_2$}
        \begin{enumerate}
          \item \textbf{app\textsubscript{1}}: evaluate $e_1$ to closure $\langle\delta^c, e^c\rangle$
          \item \textbf{app\textsubscript{2}}: evaluate $e_2$ to argument $v_2$
          \item \textbf{app\textsubscript{3}}: switch to $\delta^c$, evaluate body $e^c$;\\\emph{restore caller's environment} on return
        \end{enumerate}
      \end{block}
  \end{columns}
\end{frame}

\begin{frame}{The Open Problem: Small-Step Equivalence for STLC}
  \begin{columns}[T]
    \column{0.5\textwidth}
      \begin{block}{The Mismatch}
        \textbf{Small-step}: substitution
        \[ (\lambda\, e_1) \cdot v_2 \;\longrightarrow\; e_1[v_2/0] \]
        \textbf{Big-step}: environment lookup
        \[ (\delta, v_2) \vdash e_1 \Downarrow v \]
      \end{block}
      \medskip
      To relate them requires embedding values back as terms:\\[4pt]
      $\mathbf{embed} : \vDash T \to \emptyset \vdash T$\\[4pt]
      But closures $\langle \delta^c, e^c \rangle$ are defined in context $\Gamma^c \ne \emptyset$ and the captured environment cannot be embedded.
    \column{0.46\textwidth}
      \alert{Left open in this thesis.}
      \bigskip

      \textbf{Possible approaches:}
      \begin{itemize}
        \item Kripke logical relations
        \item Named-variable representation with explicit substitution lemmas
        \item Unified representation where both semantics share the same notion of value
      \end{itemize}
      \bigskip
      {\footnotesize An alternative formulation (``STLC-Alt'') using a \texttt{quote} constructor was explored but does not resolve the core issue.}
  \end{columns}
\end{frame}

%%%%%%%%%%%%%%%%%%%%%%%%%%%%%%%%%%%%%%%%%%%%%%%%%%%%%%%%%%%%%%%%%%%%%%%%%%%
\section{Related Work \& Conclusion}
%%%%%%%%%%%%%%%%%%%%%%%%%%%%%%%%%%%%%%%%%%%%%%%%%%%%%%%%%%%%%%%%%%%%%%%%%%%

\begin{frame}[coloredbackground,stackedUbackground]{}
  \centering\vfill
  {\usebeamerfont{title}\usebeamercolor[fg]{title}Part VI\par}
  \bigskip
  {\Large\bfseries Related Work \& Conclusion}
  \vfill
\end{frame}

\begin{frame}{Related Work}
  \begin{description}\setlength{\itemsep}{8pt}
    \item[Hutton et al.\ \textup{(``Easy as 1,2,3'')}]
      Evaluator $\to$ CPS $\to$ defunctionalization $\to$ machine.\\[2pt]
      {\footnotesize \emph{A frame in this thesis $\equiv$ a defunctionalized continuation.} Correctness by construction; different derivation route.}
    \item[Hutton \& Wright \textup{(``Calculating Exceptional Machines'')}]
      Same CPS approach extended with exceptions.\\[2pt]
      {\footnotesize The continuation type is modified to account for exception-raising vs.\ returning.}
    \item[Hutton et al.\ \textup{(``Compiling Exceptions Correctly'')}]
      Verified compiler to a stack-based machine; handlers pushed as code.\\[2pt]
      {\footnotesize Exceptions unwind until a matching handler is found.}
    \item[Hutton et al.\ \textup{(``Calculating Dependently Typed Compilers'')}]
      Agda formalization --- \emph{without} explicit correctness proofs.
  \end{description}
  \medskip
  \begin{block}{This Thesis}
    Derives machines from the \emph{structure of big-step derivations} directly,\\with \textbf{machine-checked} correctness and completeness proofs.
  \end{block}
\end{frame}

\begin{frame}{Summary}
  \begin{columns}[T]
    \column{0.5\textwidth}
      \begin{block}{Three Completed Case Studies}
        \begin{tabular}{ll}
          Arithmetic  & full chain \\
          Exceptions  & full chain + unwind mode \\
          STLC        & big-step $\leftrightarrow$ machine \\
        \end{tabular}
      \end{block}
      \medskip
      \textbf{For each language:}
      \begin{itemize}
        \item Intrinsically typed syntax
        \item Big-step semantics (total, deterministic)
        \item Abstract machine derived from frames
        \item Machine-checked completeness + correctness
        \item Denotational semantics + equivalence
      \end{itemize}
    \column{0.46\textwidth}
      \textbf{Semantic equivalence chain:}
      \[
        \text{Denotational} \leftrightarrow \text{Big-Step} \leftrightarrow \text{Machine}
      \]
      \medskip
      \textbf{Uniform proof methodology:}
      \begin{itemize}
        \item Completeness: induction on big-step derivation
        \item Correctness: totality + completeness + confluence
        \item Same pattern across all three case studies
      \end{itemize}
      \medskip
      \alert{All Agda code type-checks.}\\[3pt]
      {\footnotesize Correctness guaranteed by construction, not by testing.}
  \end{columns}
\end{frame}

\begin{frame}{Limitations \& Future Work}
  \begin{columns}[T]
    \column{0.5\textwidth}
      \textbf{Current limitations:}
      \begin{enumerate}\setlength{\itemsep}{6pt}
        \item \textbf{Totality assumption}\\[2pt]
              Big-step semantics cannot represent non-terminating programs.
        \item \textbf{Termination pragma in STLC}\\[2pt]
              The totality proof is accepted on trust, not verified structurally.
        \item \textbf{Small-step $\leftrightarrow$ big-step for STLC}\\[2pt]
              Substitution vs.\ environment lookup remains unproven.
      \end{enumerate}
    \column{0.46\textwidth}
      \textbf{Future directions:}
      \begin{itemize}\setlength{\itemsep}{6pt}
        \item Coinductive big-step or pivot to small-step as the primary specification
        \item Find a structurally terminating proof for STLC totality
        \item Prove STLC small-step equivalence via Kripke logical relations
        \item Abstract the derivation methodology into a reusable framework
      \end{itemize}
  \end{columns}
\end{frame}

\begin{frame}[coloredbackground,stackedUbackground]{}
  \centering\vfill
  {\usebeamerfont{title}\usebeamercolor[fg]{title}\Large Thank You\par}
  \bigskip
  Questions?
  \vfill
\end{frame}

\end{document}
